% GENERATED FILE. Edit canonical lessons, then run npm run generate:books.
% canonical-source-hash: fnv1a64:8013111c912f1e74
% canonical-lessons: FR-C34-il-y-a, FR-C34-cest, FR-C34-ici-la, FR-C34-voici, FR-C34-ce-cette, FR-R34-voici-tout

\chapter{Five Ways to Point}
\label{ch:fr-pointing}
\hlchaptermodality{38}

\begin{chapteropening}
\textbf{By the end of this chapter:} \emph{I can say that something exists, say what it is, hand somebody over, point at a noun, and answer the question where.}

Everything the book owns is pointed at five ways: it exists, it is a thing, it is handed over, it is pointed at in front of its noun, and it is somewhere. One Latin word runs under four of the five -- ecce, 'behold', which fused into ce, ici, voici and voila and then vanished inside all of them.
\end{chapteropening}

\section[il y a]{il y a — \textquotedblleft{}there is,\textquotedblright{} built out of \textquotedblleft{}it has there\textquotedblright{}}

\label{lesson:FR-C34-il-y-a}



\begin{quote}
Say \textbf{il a} and \textbf{il est}. You have both verbs whole, and you still cannot say that something exists.

\begin{itemize}
\raggedright
  \item \emph{Recall:} two chapters back, the three ways to ask a question — voice, \emph{est-ce que}, inversion
\end{itemize}
\end{quote}

\begin{cousinweb}
\begin{quote}
\textbf{il y a} — there is, there are.
\end{quote}

Said \emph{eel-ya}, three words run into two syllables.

\begin{quote}
\textbf{Il y a du pain.} — There is some bread.
\end{quote}

\begin{quote}
\textbf{Il y a un café.} — There is a café.
\end{quote}

Word for word it says \textbf{it has there}. The \textbf{il} stands for nobody — the same empty \emph{il} as in \emph{il est trois heures} — the \textbf{a} is \emph{avoir}, and the \textbf{y} is the piece worth stopping on.

That \textbf{y} is Latin \textbf{ibi}, \textquotedblleft{}there,\textquotedblright{} and it survives in French almost nowhere else. Spanish welded the same \emph{ibi} onto its own \emph{ha} and got \textbf{hay}; French kept the two apart and writes them still. Two languages, one Latin adverb, one idea — and only French leaves the seam visible.

One shape covers both numbers. \textbf{Il y a un café} and \textbf{il y a deux cafés} use the identical phrase, so there is no \emph{there are} to learn separately.
\end{cousinweb}

\subsection*{Your turn}
Say these aloud:

\begin{itemize}
\raggedright
  \item \emph{il y a} — \emph{eel-ya}
  \item with the partitive you already own — \textquotedblleft{}Il y a du pain, du vin, de l'eau.\textquotedblright{}
  \item with a count — \textquotedblleft{}Il y a un café\textquotedblright{} and \textquotedblleft{}Il y a deux cafés\textquotedblright{}, the same phrase twice
  \item the three pieces, slowly — \emph{il … y … a}
  \item with the chain from the last chapter, the second verb unchanged — \textquotedblleft{}Je veux du pain\textquotedblright{} and \textquotedblleft{}Il va y avoir du pain\textquotedblright{}
\end{itemize}

\begin{tcolorbox}[breakable,colback=teal!4,colframe=teal!35!black,title={Before you move on}]
What does \emph{il y a} say word for word? (\textbf{It has there.}) What is the \textbf{y}? (\textbf{Latin \emph{ibi}, \textquotedblleft{}there\textquotedblright{}} — the same word inside Spanish \emph{hay}.) Does it change for plural? (\textbf{No.}) Say \textquotedblleft{}there is some bread.\textquotedblright{}
\end{tcolorbox}
\section[c'est]{c'est — \textquotedblleft{}this is,\textquotedblright{} and the pointing phrase inside it}

\label{lesson:FR-C34-cest}



\begin{quote}
\emph{Il y a} says a thing exists. It cannot say \textbf{what} the thing is. Say \textbf{il est} and \textbf{l'eau}, and watch the apostrophe do its usual work.
\end{quote}

\begin{cousinweb}
\begin{quote}
\textbf{c'est} — this is, that is. · \textbf{ce sont} — these are, those are.
\end{quote}

\textbf{ce} is the pointing word, \textbf{est} is \emph{être}, and the \textbf{e} of \emph{ce} falls out in front of a vowel exactly as \emph{le} becomes \emph{l'}:

\begin{quote}
\textbf{C'est du café.} — That's coffee. · \textbf{C'est mon frère.} — That's my brother.
\end{quote}

For a plural French switches the verb and stops eliding, because \emph{sont} opens with a consonant:

\begin{quote}
\textbf{Ce sont des cafés.} — Those are cafés.
\end{quote}

\textbf{ce} is Latin \textbf{ecce hoc} — \emph{behold this} — a whole phrase worn down to two letters over fifteen centuries. The \emph{ecce} half is still doing work elsewhere in this language, and the next two lessons will show where.
\end{cousinweb}

\subsection*{Your turn}
Say these aloud:

\begin{itemize}
\raggedright
  \item \emph{c'est} — one syllable, \emph{say}
  \item the table — \textquotedblleft{}C'est du pain. C'est du fromage. C'est du café.\textquotedblright{}
  \item people — \textquotedblleft{}C'est mon frère. C'est ma sœur.\textquotedblright{}
  \item the plural — \textquotedblleft{}Ce sont des cafés.\textquotedblright{}
  \item the pair — \textquotedblleft{}Il y a un café\textquotedblright{} against \textquotedblleft{}C'est un café\textquotedblright{}
\end{itemize}

\begin{tcolorbox}[breakable,colback=teal!4,colframe=teal!35!black,title={Before you move on}]
What two words are inside \emph{c'est}? (\emph{\textbf{Ce}} and \emph{\textbf{est}}.) What happens for a plural? (\emph{\textbf{Ce sont}}, with no elision.) What Latin phrase is \emph{ce}? (\emph{\textbf{Ecce hoc}} — \textquotedblleft{}behold this\textquotedblright{}.) What does \emph{il y a} do that \emph{c'est} does not? (\textbf{Say a thing exists without saying what it is.})
\end{tcolorbox}
\section[ici, là, là-bas]{ici, là, là-bas — three distances}

\label{lesson:FR-C34-ici-la}



\begin{quote}
You can ask \textbf{où ?}. You have had no way at all to answer it.
\end{quote}

\begin{cousinweb}
\begin{quote}
\textbf{ici} — here. · \textbf{là} — there. · \textbf{là-bas} — over there, further off.
\end{quote}

\begin{quote}
\textbf{Où est le café ? — Ici.}
\end{quote}

\textbf{ici} is Latin \textbf{ecce hic}, \textquotedblleft{}behold here\textquotedblright{} — the same \textbf{ecce}, \textquotedblleft{}behold,\textquotedblright{} that is buried inside \textbf{c'est}. French said its pointing word so often that it fused into several separate everyday words and then disappeared inside each of them; two more of them arrive in the next lesson.

\textbf{là-bas} is \textbf{là} plus \textbf{bas}, \textquotedblleft{}low\textquotedblright{} — \emph{there below}. The picture is of something further down a road, and it marks the far end of a three-step scale English can only manage with two words and a gesture.

The accent on \textbf{là} is doing real work: it separates the adverb from \textbf{la}, the article. Same letter, same sound, and the grave mark is the only difference on the page.
\end{cousinweb}

\subsection*{Your turn}
Say these aloud:

\begin{itemize}
\raggedright
  \item the three, near to far — \textquotedblleft{}ici, là, là-bas\textquotedblright{}
  \item the question and its answer — \textquotedblleft{}Où est le café ? — Ici.\textquotedblright{}
  \item \textquotedblleft{}Le pain est là.\textquotedblright{}
  \item \textquotedblleft{}Il y a un café là-bas.\textquotedblright{}
  \item the pair the accent separates — \textquotedblleft{}la\textquotedblright{} the article, \textquotedblleft{}là\textquotedblright{} the place
  \item with what it is — \textquotedblleft{}C'est le café. Il est là.\textquotedblright{}
\end{itemize}

\begin{tcolorbox}[breakable,colback=teal!4,colframe=teal!35!black,title={Before you move on}]
Say the three distances. (\emph{\textbf{Ici, là, là-bas}}.) What is \emph{là-bas} made of? (\emph{\textbf{Là}} plus \emph{\textbf{bas}}, \textquotedblleft{}low\textquotedblright{} — \emph{there below}.) What does the accent on \emph{là} separate it from? (\textbf{The article \emph{la}}.) Answer \emph{où est le pain ?} three ways.
\end{tcolorbox}
\section[voici]{voici — \textquotedblleft{}here is,\textquotedblright{} and introducing somebody else}

\label{lesson:FR-C34-voici}



\begin{quote}
You can give your own name, ask for somebody else's, and shake hands afterwards. Say all three. Then notice that you have never been able to introduce a third person.
\end{quote}

\begin{cousinweb}
\begin{quote}
\textbf{voici} — here is, here are. · \textbf{voilà} — there is, there are.
\end{quote}

\begin{quote}
\textbf{Voici ma mère.} — This is my mother.
\end{quote}

\begin{quote}
\textbf{Voilà mon père.} — That's my father.
\end{quote}

Both are frozen commands. \textbf{voici} is \textbf{vois} + \textbf{ci}, \textquotedblleft{}see here\textquotedblright{}; \textbf{voilà} is \textbf{vois} + \textbf{là}, \textquotedblleft{}see there\textquotedblright{}. Nobody hears an order in them any more, and the verb inside has not been conjugated for centuries.

The \textbf{ci} and the \textbf{là} are the two words you took last lesson, and the \emph{vois} in front of them carries the same Latin \textbf{ecce}, \textquotedblleft{}behold,\textquotedblright{} that is buried inside \emph{ce} and \emph{ici}. So French points four ways off one ancient root — \emph{ce}, \emph{ici}, \emph{voici}, \emph{voilà} — and each of the four froze at a different moment.

Between them and \emph{c'est}, you can now hand somebody over:

\begin{quote}
\textbf{Voici ma sœur. — Enchanté.}
\end{quote}
\end{cousinweb}

\subsection*{Your turn}
Say these aloud:

\begin{itemize}
\raggedright
  \item \emph{voici} and \emph{voilà}
  \item the household — \textquotedblleft{}Voici ma mère. Voici mon père.\textquotedblright{}
  \item \textquotedblleft{}Voici mon frère. Voici ma sœur.\textquotedblright{}
  \item \textquotedblleft{}Voici mon ami.\textquotedblright{} then the reply, \textquotedblleft{}Enchanté.\textquotedblright{}
  \item the four pointing words in one breath — \textquotedblleft{}ce, ici, voici, voilà\textquotedblright{}
  \item your own name, then somebody else's — \textquotedblleft{}Je m'appelle…\textquotedblright{}, \textquotedblleft{}Voici…\textquotedblright{}
\end{itemize}

\begin{tcolorbox}[breakable,colback=teal!4,colframe=teal!35!black,title={Before you move on}]
What two words are frozen inside \emph{voici}? (\emph{\textbf{Vois}} and \emph{\textbf{ci}} — \textquotedblleft{}see here\textquotedblright{}.) And inside \emph{voilà}? (\emph{\textbf{Vois}} and \emph{\textbf{là}}.) Introduce your mother, and answer yourself. (\textbf{Voici ma mère. — Enchanté.})
\end{tcolorbox}
\section[ce, cet, cette, ces]{ce, cet, cette, ces — pointing at a noun}

\label{lesson:FR-C34-ce-cette}



\begin{quote}
Say \textbf{le}, \textbf{la}, \textbf{les}. Those say \emph{the}. Nothing you own says \emph{this} in front of a noun — \emph{c'est} and \emph{voici} both need a whole sentence.
\end{quote}

\begin{cousinweb}
\begin{quote}
\textbf{ce}, \textbf{cet}, \textbf{cette}, \textbf{ces} — this, that, these, those, standing in front of a noun.
\end{quote}

The choice follows the article you already know, one for one:

\noindent
\begin{tabularx}{\linewidth}{@{}>{\raggedright\arraybackslash}X>{\raggedright\arraybackslash}X@{}}
\toprule
\textbf{the noun takes} & \textbf{you say} \\
\midrule
\emph{le} & \textbf{ce} pain \\
\emph{le}, before a vowel & \textbf{cet} ami \\
\emph{la} & \textbf{cette} mère \\
\emph{les} & \textbf{ces} cafés \\
\bottomrule
\end{tabularx}

Only the second line is new machinery, and it is machinery you have already met from the other side. \emph{Le ami} is impossible, so French writes \textbf{l'ami}. \emph{Ce ami} is impossible for the identical reason — and here French borrows a \textbf{t} instead of eliding, giving \textbf{cet ami}.

The result is a small trap for the ear: \textbf{cet ami} and \textbf{cette amie} are pronounced exactly alike, and only the sentence around them says which was meant.

All four are Latin \textbf{ecce iste}, \textquotedblleft{}behold that one\textquotedblright{} — the same \textbf{ecce} inside \emph{ce}, \emph{voici} and \emph{voilà}, welded this time to a different second half.
\end{cousinweb}

\subsection*{Your turn}
Say these aloud:

\begin{itemize}
\raggedright
  \item the article first, then the pointer — \textquotedblleft{}le pain, ce pain\textquotedblright{}
  \item \textquotedblleft{}la mère, cette mère\textquotedblright{}
  \item \textquotedblleft{}les cafés, ces cafés\textquotedblright{}
  \item the borrowed \emph{t} — \textquotedblleft{}l'ami, cet ami\textquotedblright{}
  \item the animals — \textquotedblleft{}ce chien, ce chat\textquotedblright{}
  \item the two that sound alike — \textquotedblleft{}cet ami\textquotedblright{}, \textquotedblleft{}cette amie\textquotedblright{}
\end{itemize}

\begin{tcolorbox}[breakable,colback=teal!4,colframe=teal!35!black,title={Before you move on}]
Which shape goes in front of a \emph{le}-word? (\emph{\textbf{Ce}}.) A \emph{la}-word? (\emph{\textbf{Cette}}.) A plural? (\emph{\textbf{Ces}}.) Why does \emph{cet} exist? (\textbf{To keep two vowels apart}, the same job the apostrophe does in \emph{l'ami}.) Which two of the four sound identical? (\emph{\textbf{Cet}} and \emph{\textbf{cette}}.)
\end{tcolorbox}
\section[Practice]{Voici tout — the chapter, run cold}

\label{lesson:FR-R34-voici-tout}



\begin{quote}
One Latin word is inside four things you learned this chapter. Name it, then name the four.
\end{quote}

\subsection*{Your turn: the table, pointed at five ways}
Everything the book has put on a table, pointed at five different ways. Say each line before you read the next.

Say these aloud:

\begin{itemize}
\raggedright
  \item it exists — \textquotedblleft{}Il y a du pain, du vin, de l'eau.\textquotedblright{}
  \item what it is — \textquotedblleft{}C'est du fromage. C'est du café. C'est du lait.\textquotedblright{}
  \item handed over — \textquotedblleft{}Voici le thé. Voici le sel. Voici le sucre.\textquotedblright{}
  \item pointed at — \textquotedblleft{}ce beurre, cet œuf, cette eau, ces cafés\textquotedblright{}
  \item the loan and the inheritance side by side — \emph{café} and \emph{thé} came from far away; \emph{pain}, \emph{vin} and \emph{lait} never left
\end{itemize}

\subsection*{Your turn: the distant band — people, animals, the body}
Now reach back. Every person, every animal and every part of the body this book has taught, introduced or pointed at.

Say these aloud:

\begin{itemize}
\raggedright
  \item the household, handed over — \textquotedblleft{}Voici ma mère, mon père, mon frère, ma sœur\textquotedblright{}
  \item the four Latin cousins those keep — \emph{mater, pater, frater, soror}
  \item \textquotedblleft{}Voici mon ami. Voici l'enfant. Voici la famille.\textquotedblright{} then \textquotedblleft{}Enchanté.\textquotedblright{}
  \item the two animals, pointed at — \textquotedblleft{}ce chien, ce chat\textquotedblright{}
  \item the body — \textquotedblleft{}cette tête, cette main, ce nez, cet œil, cette bouche, ce ventre\textquotedblright{}
  \item the one that needed a borrowed \emph{t}, and why — \textquotedblleft{}cet œil\textquotedblright{}
\end{itemize}

\subsection*{Your turn: where}
\begin{itemize}
\raggedright
  \item \emph{Say it:} \textquotedblleft{}Où est le café ?\textquotedblright{} then all three answers — \textquotedblleft{}Ici. Là. Là-bas.\textquotedblright{}
  \item \emph{Say it:} \textquotedblleft{}Il y a un café là-bas.\textquotedblright{}
  \item \emph{Ask:} the same question the other two ways you know
\end{itemize}

\begin{tcolorbox}[breakable,colback=teal!4,colframe=teal!35!black,title={Before you move on}]
Which phrase says a thing exists without saying what it is? (\emph{\textbf{Il y a}}.) Which says what it is? (\emph{\textbf{C'est}}.) Which hands it to somebody? (\emph{\textbf{Voici}}.) Which points at it in front of its noun? (\emph{\textbf{Ce, cet, cette, ces}}.) And what single Latin word is buried in \emph{ce}, \emph{ici}, \emph{voici} and \emph{voilà}? (\emph{\textbf{Ecce}} — \textquotedblleft{}behold\textquotedblright{}.)
\end{tcolorbox}
